Please follow the NeurIPS template https://nips.cc/Conferences/2020/PaperInformation/StyleFiles (Links to an external site.) for the project report. It's OK to use Microsoft Word. But please make sure it look as closely possible to the LaTex template as possible. Significant difference will lead to point deductions.

Your final report is expected to be 6-8 pages excluding references, in accordance with the length requirements for a NeurIPS paper. It should have exactly the following format:

    Introduction: problem definition and motivation
    Related Work: literature survey (at least 6 papers)
    Methods – Overview of your method – Intuition on why should it be better than the baseline – Details of models and algorithms that you developed
    Experiments – Description of your testbed and a list of questions your experiments are designed to answer – Details of the experiments and results
    Conclusion: discussion and future work

Failing to follow such a format will lead to point deductions.

The grading breakdown for the final report is as follows:

    10% for introduction and literature survey
    30% for proposed method (soundness and originality)
    30% for correctness, completeness, and difficulty of experiments and figures
    10% for empirical analysis of results and methods
    20% for quality of writing (clarity, organization, flow, etc.)

If you use any open-source code, please disclose it in the end of the introduction section. Failing to do so will lead to point deductions.

Since there is a strict deadline ahead to submit the final grades to the university, no late days are allowed to be used for the final report, as has been mentioned in the class.

We will select 5 projects (at most) that we found the most innovative (technically novel, a new task setting, new empirical insights, etc). Members of the selected projects will receive 5 extra points toward your final grade.